\subsection{箭图的态射}

设 C 与 C' 为箭图。从 C 到 C' 的箭图态射是一个偶 \((u,v)\) ，其中 \(u\colon\operatorname{Som}(\mathrm{C})\to\operatorname{Som}(\mathrm{C}')\) 与 \(v\colon\operatorname{Fl}(\mathrm{C})\to\operatorname{Fl}(\mathrm{C}')\) 是满足 \(o_{C'}\circ v=u\circ o_{C}\) 与 \(t_{C'}\circ v=u\circ t_{C}\) 的映射。

设 \(\varphi = (u, v)\) 为从 C 到 C' 的箭图态射。映射 \(u\) 与 \(v\) 分别记作 \(\operatorname{Som}(\varphi)\) 与 \(\operatorname{Fl}(\varphi)\) 。若 \(a\) 是 C 的顶点，则滥用记号，把 \(\varphi(a)\) 记为顶点 \(a\) 在 \(\varphi\) 下的象，我们将说 \(\varphi(a)\)

是 \(a\) 在 \(\varphi\) 下的象。类似地，若 \(f\) 是 C 的箭，则滥用记号，把 \(\varphi(f)\) 记为 C' 的箭 \(v(f)\) ，我们将说它是 \(f\) 在 \(\varphi\) 下的象。

设 C、\(C'\) 、\(C''\) 为三个箭图，\(\varphi = (u, v)\) 为从 C 到 \(C'\) 的态射，\(\varphi' = (u', v')\) 为从 \(C'\) 到 \(C''\) 的态射。则 \((u' \circ u, v' \circ v)\) 是从 C 到 \(C''\) 的态射，记作 \(\varphi' \circ \varphi\) ，称为 \(\varphi'\) 与 \(\varphi\) 的复合。

我们把 \(\mathrm{Id}_{\mathrm{C}}\) 记为从 C 到自身的态射 \((\mathrm{Id}_{\mathrm{Som}(\mathrm{C})}, \mathrm{Id}_{\mathrm{Fl}(\mathrm{C})})\) 。

设 \(\varphi\) 为从 C 到 C' 的箭图态射。为使 \(\varphi\) 是同构，必须且只需映射 \(\operatorname{Som}(\varphi)\) 与 \(\operatorname{Fl}(\varphi)\) 都是双射（参见 E, IV, 第 6 页）。

这样，若把集合 S 与 F 上的箭图结构称为两个映射 \(o, t: \mathrm{F} \to \mathrm{S}\) 的数据，就可以把箭图态射取作这个结构的态射（E, IV, 第 11 页）。

设 \(\varphi\) 为从 C 到 C' 的箭图态射。存在 C' 的唯一子箭图，其顶点集与箭集分别是 \(\mathrm{Som}(\varphi)\) 与 \(\mathrm{Fl}(\varphi)\) 的象。把它记作 \(\varphi(\mathrm{C})\) ，称为 C 在 \(\varphi\) 下的象。称 \(\varphi\) 是单射的，若映射 \(\mathrm{Som}(\varphi)\) 与 \(\mathrm{Fl}(\varphi)\) 都是单射；这时 \(\varphi\) 诱导一个从 C 到其象上的同构。

设 \(\varphi\) 与 \(\psi\) 为从 C 到 C' 的箭图态射。存在 C 的唯一子箭图，其顶点与箭分别是 C 中在 \(\varphi\) 与 \(\psi\) 下有相同象的顶点与箭。我们称之为 \(\varphi\) 与 \(\psi\) 的均衡子。

